\section{Auto-loan fixed-confidence policy-rule replay}
\label{sec:autoloan-llm-replay}

We evaluate the generator-augmentation mechanism on an online auto-loan replay. The experiment instantiates the finite contextual policy-class formulation. A unit context \(W\) is a loan application record, the action is an offered margin price, and each arm is a pricing rule mapping the application context to an offer. The two rules are the historical company pricing rule and a support-constrained improved rule that chooses a multiplier from \(\{0.8,0.9,1.0,1.1,1.2\}\) times the historical offer.

A pull samples an application from the later-period evaluation pool, applies one of the two pricing rules, draws acceptance from the fitted replay demand model, and returns realized margin revenue. For each LLM method, the proxy score is generated for the same application--rule pair and is converted to predicted revenue by multiplying the model's predicted acceptance probability by the offered margin price. The proxy is not used as a reward or as a direct ranking signal; it enters only through same-unit reward--proxy covariance.

Unlike a fixed-budget budget--accuracy diagnostic, the main replay uses a fixed-confidence stopping rule. In each repetition, the algorithm samples until the lower confidence bound of the currently estimated better pricing rule exceeds the upper confidence bound of the other rule. We report stopping times at confidence level \(\delta=0.05\), empirical correctness at stopping, and armwise reward--proxy correlations. The reward-only baseline uses the same stopping rule without a proxy. Each LLM method uses the lower-confidence-bound Gen-CV estimator; the oracle control-variate row is retained only as a diagnostic ceiling.

We compare four methods: reward-only EGE, local Qwen in-context scoring without fine-tuning, LoRA fine-tuned Qwen using the same base checkpoint, and GPT API in-context scoring without fine-tuning. The two Qwen variants use the same open-weight base checkpoint; the only difference is whether task-specific LoRA fine-tuning is applied. All LLM variants use the same paired scoring rows and the same fixed-confidence replay code; only the proxy predictions differ.
